\documentclass{article}

\usepackage[final]{neurips_2020}

\usepackage[utf8]{inputenc}
\usepackage[T1]{fontenc}
\usepackage{hyperref}
\usepackage{url}
\usepackage{booktabs}
\usepackage{amsfonts}
\usepackage{amsmath}
\usepackage{nicefrac}
\usepackage{microtype}
\usepackage{graphicx}
\usepackage{xcolor}
\usepackage{multirow}

% \note{...} prints TODO reminders in blue. Delete notes before final submission.
\newcommand{\note}[1]{\textcolor{blue}{{#1}}}

\title{
  Difficulty-Aware Language Routing and Cross-Lingual Self-Consistency \\
  for Korean Mathematical Reasoning \\
  \vspace{1em}
  \small{\normalfont Korea University COSE461 Final Project}
}

\author{
  이윤제 (Yunje Lee) \\
   Department of Computer Science \\
   Team \note{N} \\
   2022320317
  \And
   김상준 (Sangjun Kim) \\
   Department of Computer Science \\
   Team \note{N} \\
   2022320306
  \And
   원준서 (Junseo Won) \\
   Department of Computer Science \\
   Team \note{N} \\
   2022320302
}

\begin{document}

\maketitle

\begin{abstract}
Small language models reason markedly better in English than in Korean on
mathematical word problems. We trace part of this gap to the
distillation pipeline itself: even a strong teacher (Gemini) produces
lower-quality chain-of-thought (CoT) in Korean (85.7\% correct) than in
English (97.5\%), so naively distilling Korean CoT transfers a noisier
signal. We address this with a two-stage cross-lingual framework. At the
\textbf{data level}, our main contribution \textbf{Difficulty-Aware
Language Routing (DALR)} decides, per problem, whether to learn from a
Korean or an English teacher CoT based on teacher correctness: English
CoT is used only as a ``bridge'' on problems the Korean teacher fails.
A controlled ablation (DALR vs.\ random English placement) shows DALR's
gain stems from \emph{where} English reasoning is routed (hard problems)
rather than from added data. At the \textbf{inference level}, we apply
\textbf{Cross-Lingual Self-Consistency (XLSC)}, sampling each problem in
both Korean and English and aggregating by majority vote, with a
\textbf{Cascade} extension that resolves ties through additional English
samples. On HRM8K (Korean) with Qwen2.5-3B, DALR achieves the best
single-model Korean accuracy, and DALR\,+\,XLSC further improves
performance, with the same trend reproduced on Llama-3.2-3B and
EXAONE-3.5-2.4B.
\end{abstract}

This project report is limited to 8 content pages (excluding reference); i.e. you can submit reports with 1-8 pages. Please make your report brief and concise, so that there is no redundant contents. There is no bonus point on submitting full 8 pages report.

\section{Introduction}
Large language models perform strong multi-step mathematical reasoning,
but small models (e.g.\ 3B parameters) lag far behind, and the gap is
worse outside English. For Qwen2.5-3B-Instruct, the model solves
\textbf{71.65\%} of English problems (GSM8K) but only \textbf{53.90\%}
of Korean problems (HRM8K). Practical deployments often cannot afford
large models yet must serve non-English users, making this gap
consequential.

A standard remedy is knowledge distillation: a strong teacher generates
chain-of-thought (CoT) solutions, and a small student is fine-tuned on
them. However, in the cross-lingual setting the teacher itself is
weaker in Korean: in our data Gemini produces correct CoT for
\textbf{97.5\%} of English problems but only \textbf{85.7\%} of Korean
problems. Distilling Korean CoT therefore transfers a degraded signal,
while English continues to improve, widening rather than closing the
gap.

This motivates our central question: \emph{can the teacher's---and a
model's---stronger English reasoning be used to improve a small
student's Korean mathematical reasoning?} Our working assumption is
that mathematical \emph{logic} is largely language-agnostic, so
connecting English reasoning patterns to Korean problems should help.

We pursue this through a two-stage cross-lingual framework. At the
\textbf{data level}, our main contribution \textbf{DALR} routes each
problem to a Korean or English teacher CoT based on teacher
correctness. At the \textbf{inference level}, \textbf{XLSC} samples in
both languages and votes, while the \textbf{Cascade} extension
resolves ties by additional English samples.

\paragraph{Contributions.}
\begin{itemize}
  \item We propose \textbf{DALR}, a difficulty-aware data-construction
  method for cross-lingual CoT distillation, and show via a controlled
  ablation (random vs.\ routed English bridges) that its gains stem from
  \emph{routing}, not added data.
  \item We propose \textbf{XLSC} and its \textbf{Cascade} extension,
  inference-time aggregation methods that exploit the bilingual
  reasoning capability instilled by DALR.
  \item We report results on HRM8K and GSM8K with three small models
  (Qwen2.5-3B, Llama-3.2-3B, EXAONE-3.5-2.4B) and provide a statistical
  analysis (bootstrap confidence intervals and McNemar tests).
\end{itemize}

\section{Related Work}
\paragraph{Chain-of-thought and its distillation.}
Chain-of-thought prompting elicits step-by-step reasoning in large
language models \citep{wei2022cot}, and a line of work distills such
reasoning into smaller students by fine-tuning on teacher-generated
CoT. \note{Add 1--2 CoT-distillation citations (e.g.\ Magister 2022, Ho
2022, Hsieh 2023).}

\paragraph{Cross-lingual reasoning.}
Prior work transfers reasoning across languages via translated prompts,
cross-lingual instruction tuning, or English-pivot inference
\note{(cite MGSM/xCoT)}. Unlike methods that augment \emph{prompts} at
inference, DALR selects the training CoT \emph{language} per problem
according to teacher quality, internalizing cross-lingual signal during
fine-tuning rather than at inference.

\paragraph{Self-consistency.}
Self-consistency samples multiple chains from one model and votes on
the final answer \citep{wang2022selfconsistency}. XLSC extends this by
sampling in two languages from the same model and exploiting their
decorrelated errors; Cascade further allocates inference budget only to
unresolved (tied) cases.

\section{Approach}
\subsection{Overview}
We train a single small student with DALR-constructed data and aggregate
its predictions at inference with XLSC and Cascade. Both stages inject
cross-lingual signal but at different points: DALR at training time
(deciding the language of the CoT to learn from per problem), and
XLSC/Cascade at inference time (sampling and voting in both languages).
\note{Add Figure 1: two-stage pipeline diagram.}

\subsection{Difficulty-Aware Language Routing (DALR)}
\label{sec:dalr}
Let a training problem $i$ have a Korean teacher CoT $c_i^{\text{KO}}$
and an English teacher CoT $c_i^{\text{EN}}$, each labeled correct
($v_i^{\text{KO}}, v_i^{\text{EN}} \in \{0,1\}$) by exact-match on the
numeric answer. For each Korean training problem we choose the
training target as:
\begin{itemize}
  \item if $v_i^{\text{KO}} = 1$: train on $c_i^{\text{KO}}$
  (Korean CoT, natural same-language reasoning);
  \item else if $v_i^{\text{EN}} = 1$: train on $c_i^{\text{EN}}$ as an
  \textbf{English bridge} (Korean question, English solution);
  \item else: discard the example.
\end{itemize}
Teacher failure in Korean is used as a difficulty signal: a reliable
English solution is then a higher-quality target than the unreliable
Korean one. Applied to GSM8K-ko's 7{,}473 training problems, this yields
6{,}407 Korean and 920 English-bridge examples (7{,}327 total). We
denote the resulting model \textbf{E}.

\paragraph{Ablation: routed vs.\ random.}
To isolate the effect of routing from the effect of adding English
data, we construct \textbf{E\_random}, which uses the \emph{same} 920
English bridges placed on randomly chosen \emph{easy} (Korean-solved)
problems instead of the hard ones, with everything else identical.

\subsection{Cross-Lingual Self-Consistency (XLSC)}
\label{sec:xlsc}
At inference time, for each Korean test problem we generate $N$ samples
with a Korean system prompt and $N$ samples with an English system
prompt from the \emph{same} DALR-trained model, all with sampling
temperature $T$. We extract a numeric answer from each chain and take a
$2N$-way majority vote. Because the model has internalized both
languages through DALR, both KO and EN samples are reasonably
accurate; because they reason in different languages, their errors are
partially decorrelated.

\subsection{Cascade XLSC}
\label{sec:cascade}
A non-trivial fraction of problems are tied in the $2N$-way vote
\note{(report tie rate from experiments)}. \textbf{Cascade XLSC}
addresses ties by generating one additional English sample for tied
problems, yielding a $(2N{+}1)$-way odd vote that admits no ties. We
break ties in favor of English because tied cases in our analysis show
slightly higher accuracy for the English plurality, and because DALR's
English bridges target precisely the harder (Korean-failing) problems.

\subsection{Baselines}
We compare against: \textbf{A} (base, no fine-tuning), \textbf{B}
(English-only CoT distillation), \textbf{C} (Korean-only CoT
distillation), \textbf{D} (bilingual 50/50 mix). \textbf{E} is DALR;
\textbf{E\_random} is its routing ablation. All baselines share the
student model, LoRA configuration, and training hyperparameters.

\section{Experiments}
\subsection{Data}
\paragraph{Training.} GSM8K \note{(cite Cobbe et al.\ 2021)} provides
7{,}473 English grade-school math problems. We use GSM8K-ko, a
machine-translated Korean version. The teacher (Gemini) generates CoT
in each language; we retain only CoTs whose final numeric answer
matches the gold label (English: 7{,}283 valid; Korean: 6{,}407 valid).

\paragraph{Evaluation.} We evaluate on \textbf{HRM8K}
\note{(cite HAERAE-HUB)}, a human-reviewed Korean translation of the
GSM8K test set (1{,}319 problems), and on \textbf{GSM8K} (English,
1{,}319 problems). HRM8K is our main benchmark because its translation
quality is substantially higher than GSM8K-ko.

\subsection{Evaluation method}
We report exact-match accuracy on the final numeric answer, extracted
by a regular expression and compared to the gold answer with a small
numeric tolerance. Single-model evaluation uses greedy decoding
(temperature 0, repetition penalty 1.1) with a zero-shot CoT system
prompt in the target language. XLSC uses temperature \note{0.7} with
$N{=}\note{3}$ samples per language ($2N{=}6$-way vote); Cascade adds
one English sample for each tied problem.

We report 95\% bootstrap confidence intervals on accuracy and McNemar
tests for paired model comparisons.

\subsection{Experimental details}
The main student is \textbf{Qwen2.5-3B-Instruct}, fine-tuned with LoRA
($r{=}32$, $\alpha{=}64$, all linear modules) for 3 epochs at learning
rate $2\times 10^{-4}$, effective batch size 8 (per-device 1, grad
accumulation 8), cosine schedule, on a single RTX~5060~Ti. We
additionally train \textbf{Llama-3.2-3B-Instruct} and
\textbf{EXAONE-3.5-2.4B-Instruct} with the same recipe to assess
robustness across student architectures.

\section{Results}
\label{sec:results}

\begin{table}[t]
\centering
\small
\begin{tabular}{lccc}
\toprule
\textbf{Setup} & \textbf{HRM8K (KO)} & \textbf{95\% CI (KO)} & \textbf{GSM8K (EN)} \\
\midrule
A (Base zero-shot)           & 53.90          & [51.02, 56.48] & 71.65 \\
B (English CoT)              & 54.44          & [51.78, 57.32] & 74.22 \\
C (Korean CoT)               & 59.51          & [56.63, 62.17] & 70.74 \\
D (Bilingual mix)            & 58.98          & [56.25, 61.79] & 74.98 \\
E (DALR)                     & 61.79          & [59.29, 64.44] & 69.37 \\
E\_random (ablation)         & 58.83          & [56.10, 61.79] & 72.71 \\
\midrule
E + XLSC ($N{=}3$, $T{=}0.7$)        & 68.69          & [66.19, 71.19] & \note{--} \\
E + Cascade XLSC                     & \textbf{69.98} & [67.40, 72.48] & \note{--} \\
\bottomrule
\end{tabular}
\caption{Main results on HRM8K (Korean, 1{,}319 problems) and GSM8K
(English, 1{,}319 problems) with the Qwen2.5-3B student. E (DALR) is
the best single-model setup; \textbf{E + Cascade XLSC} attains the best
overall accuracy on HRM8K (\textbf{+8.19} pts over E alone). 95\% CI
from 1{,}000-resample bootstrap.}
\label{tab:main}
\end{table}

\begin{table}[t]
\centering
\small
\begin{tabular}{lccc}
\toprule
\textbf{Student model} & \textbf{C (KO CoT)} & \textbf{E (DALR)} & \textbf{E + XLSC} \\
\midrule
Qwen2.5-3B        & 59.51 & 61.79 & \textbf{68.69} \\
Llama-3.2-3B      & \note{TBD}  & \note{TBD}  & \note{TBD} \\
EXAONE-3.5-2.4B   & \note{TBD}  & \note{TBD}  & \note{TBD} \\
\bottomrule
\end{tabular}
\caption{Robustness across student models on HRM8K. \note{Fill from
final-eval results.}}
\label{tab:robustness}
\end{table}

\paragraph{DALR improves Korean, and routing (not data) is the cause.}
On HRM8K, E (61.79) exceeds the Korean-only baseline C (59.51) by 2.28
points and the bilingual mix D (58.98) by 2.81 points, becoming the
best single model. The ablation E\_random uses the \emph{same} 920
English bridges placed on easy problems and drops to 58.83 (below C),
a 2.96-point gap from E (McNemar $p{=}0.030$). This isolates
\emph{routing}---rather than added data---as the source of DALR's
gain.

\paragraph{A Korean--English trade-off.}
DALR raises Korean accuracy but slightly lowers English (E: 69.37 vs.\
C: 70.74). This is consistent with English bridges being concentrated
on \emph{hard} Korean problems, which biases the model toward Korean
question contexts even when producing English solutions. The trade-off
motivates inference-time aggregation across both languages.

\paragraph{XLSC and Cascade further improve Korean accuracy.}
Applying XLSC ($N{=}3$ Korean + $N{=}3$ English samples, $T{=}0.7$) to
the same E checkpoint lifts HRM8K from 61.79 to \textbf{68.69}
(\textbf{+6.90} pts; McNemar $p{<}0.001$ vs.\ E). Of the 1{,}319
problems, 208 (15.8\%) end up tied in the 6-way vote; Cascade XLSC
resolves them with one extra English sample and gains another
\textbf{1.29} pts (\textbf{69.98}; McNemar $p{=}0.003$ vs.\ XLSC), with
23 ties flipping to correct and only 6 flipping to wrong. The full
pipeline (E + Cascade) improves over E by \textbf{8.19} pts
($p{<}0.001$).

\paragraph{Robustness across students.}
The same qualitative pattern---DALR matches or exceeds C on Korean,
XLSC further improves over DALR---holds on Llama-3.2-3B and
EXAONE-3.5-2.4B \note{(fill numbers in Table~\ref{tab:robustness})},
suggesting the framework is not specific to Qwen.

\section{Analysis}
\paragraph{Why does routing matter?}
Both E and E\_random use the same 920 English CoTs and the same total
training set size; they differ only in which 920 Korean training items
are \emph{replaced} by English CoT. E replaces the items the teacher
got wrong in Korean (high-difficulty), while E\_random replaces
randomly chosen items the teacher got right. E\_random therefore
discards reliable Korean training signal and substitutes mismatched
English CoT, while E substitutes only on items where the Korean signal
was unreliable in the first place. This explains why E\_random falls
below C, and why E exceeds it.

\paragraph{Where does DALR help?}
\note{Add a per-difficulty breakdown: stratify HRM8K problems by how
many of the baselines (A,B,C,D) solve them, and report E's accuracy
per stratum. Expected: largest gains on medium-difficulty problems
where Korean-only reasoning is unreliable but English bridges help.}

\paragraph{XLSC and Cascade.}
\note{Report the fraction of HRM8K problems whose KO majority answer
disagrees with the EN majority answer; report the fraction of tied
problems and the accuracy on tied vs.\ non-tied subsets; show that
Cascade improves accuracy primarily on the tied subset, justifying its
targeted budget allocation.}

\paragraph{Statistical robustness.}
All accuracies in Table~\ref{tab:main} are reported with 95\%
percentile bootstrap intervals (1{,}000 resamples). Paired McNemar
tests on HRM8K confirm the main claims:
E vs.\ E\_random $p{=}0.030$;
E vs.\ D $p{=}0.045$;
E + XLSC vs.\ E $p{<}0.001$;
Cascade vs.\ XLSC $p{=}0.003$;
Cascade vs.\ E $p{<}0.001$.
On GSM8K, E shows a significant drop relative to D ($p{<}0.001$) and
E\_random ($p{=}0.016$), confirming the deliberate Korean--English
trade-off described above.

\section{Conclusion}
We studied how to improve Korean mathematical reasoning in small
language models by leveraging stronger English reasoning at two
points in the pipeline. Our main contribution \textbf{DALR} routes
teacher CoT language per problem according to teacher correctness; a
controlled ablation shows the benefit comes from \emph{routing}
rather than added data. \textbf{XLSC} and its \textbf{Cascade}
extension exploit the bilingual capability instilled by DALR at
inference time, further improving Korean accuracy without retraining.
\textbf{Limitations.} HRM8K is a high-quality translation of GSM8K
rather than natively Korean math, so gains may partly reflect
translated-problem solving; the DALR-vs-Korean-only margin is modest
in absolute terms; we use a single teacher (Gemini) and three 3B-class
students. Future work includes natively Korean math benchmarks, larger
students, and integrating DALR's quality signal into peer
distillation.

After you are done writing this report, save it as a .pdf file, and submit through BlackBoard. You don't have to submit a copy per person; submitting once per team is fine.

\bibliographystyle{unsrt}
\bibliography{references}

\newpage
\appendix

\section*{Appendix: Team contributions}
\begin{itemize}
  \item \textbf{이윤제 (2022320317)}: \note{e.g., DALR design and
  implementation, E/E\_random training and evaluation, statistical
  analysis, paper writing (Sections~\ref{sec:dalr} and~\ref{sec:results}).}
  \item \textbf{김상준 (2022320306)}: \note{e.g., XLSC/Cascade design
  and implementation, tie analysis, paper writing
  (Sections~\ref{sec:xlsc} and~\ref{sec:cascade}).}
  \item \textbf{원준서 (2022320302)}: \note{e.g., baseline training
  (A--D), robustness experiments on Llama and EXAONE, paper writing
  (Experiments section).}
\end{itemize}

\end{document}
